\section{Dataset, software and computing}
\label{sec:data}

\subsection{Dataset}
The STAR file catalog was queried with \fcmd{get_file_list.pl} for
\fcmd{filetype=daq_reco_picoDst, trgsetupname=production_19GeV_2019}. Three productions exist:
P21ic (SL21c), P23id (SL23d) and P24iy (SL24y).

\choice{Use P24iy, the most recent production of the Run19 19.6\,GeV data (library SL24y).
The picoDst format is the one defined in \fpath{StRoot/StPicoEvent} of
\url{https://github.com/star-bnl/star-sw}; the SL24y copy of that package is used to read it.}

For P24iy the catalog lists 100\,560 files on HPSS, 99\,273 on \texttt{local} storage
($1.42\times10^9$ events) and 100\,958 NFS entries. For the latest runs, the NFS paths checked
(first 3000 entries ordered by run number) do not exist on disk.

\choice{Read the \texttt{local} replicas through XRootD,
\fpath{root://xrdstar.rcf.bnl.gov:1095//home/starlib/...} (absolute path, double slash).
The full list is \fpath{filelists/P24iy_19GeV_all_xrd.list}; the catalog dump is
\fpath{filelists/local_P24iy_all.txt}.}

\paragraph{Test sample.} One file from the last run in the catalog, run 20093036:
\inputFile{} (\nChainEvents{} events in the TTree). The first \nProcEvents{} events were processed
(\fpath{filelists/test.list}).

\caveat{A single file of a single run is not representative of the whole dataset (beam conditions,
detector status, luminosity vary run by run). Statistics are adequate for shapes and mean values,
not for run-dependent or tray-dependent conclusions.}

\subsection{Software and computing}
\label{sec:software}
\begin{table}[h]
\centering
\caption{Software environment, from \fpath{logs/provenance.txt} (version banners on the node) and the run log.}
\small
\begin{tabular}{lp{12.5cm}}
\toprule
Node & starsub01.sdcc.bnl.gov, \hostOS \\
Container & \containerImage{} (\containerOS) \\
STAR library & SL24y (\fpath{/star/nfs4/AFS/star/packages/SL24y}) \\
ROOT & \rootVersion{} (STAR spack build) \\
Compiler & gcc \gccVersion \\
Executable & \fcmd{root4star} (SL24y, .sl73\_x8664\_gcc485) \\
Run date & \runDate{} (EST) \\
\bottomrule
\end{tabular}
\end{table}

\choice{The interactive nodes run AlmaLinux 9 while the STAR libraries are SL7 builds. All builds
(\fcmd{cons}) and jobs run in the SL7 container via the wrapper \fpath{scripts/sl7.sh}
(\fcmd{apptainer exec -B /gpfs01,/gpfs,/star,/cvmfs,/direct,/sdcc ... tcsh -c "starver SL24y; ..."}).}

Working directory: \fpath{/gpfs01/star/pwg/dongx/tof/TOFME}. Code:
\begin{itemize}
\item \fpath{StRoot/StTofMEQAMaker/StTofMEQAMaker.{h,cxx}} --- QA maker (\fcmd{StMaker}) behind \fcmd{StPicoDstMaker} in read mode.
\item \fpath{macros/runTofMEQA.C} --- chain macro; sets cuts and trigger IDs.
\item \fpath{scripts/run_qa.sh N list out log} --- runs the macro in the container.
\item \fpath{macros/drawTofMEQA.C}, \fpath{scripts/draw_qa.sh} --- quick-look plots on the node.
\item \fpath{StRoot/StTofMEMaker}, \fpath{macros/runTofME.C}, \fpath{macros/makeGeom.C}, \fpath{scripts/run_me.sh},
  \fpath{scripts/me_test.sh}, \fpath{scripts/select_runs.py}, \fpath{scripts/me_v3.sub}, \fpath{scripts/merge_v3.sh}
  --- projection, mixed-event and $m^2$ step (Sec.~\ref{sec:me}--\ref{sec:m2}).
\item Note figures/tables: \fpath{analysis_note/scripts/make_figures.C}, \fpath{make_numbers.C}, \fpath{make_me.C}, \fpath{make_me_runs.C}, \fpath{make_m2.C}, \fpath{parse_log.py}, run by \fpath{build.sh} on the local copy \fpath{analysis_note/data/qa_test_10k.root}.
\end{itemize}

\choice{Only the picoDst branches \fcmd{Event}, \fcmd{Track}, \fcmd{BTofHit} and \fcmd{BTofPidTraits}
are enabled (\fcmd{StPicoDstMaker::SetStatus}), since nothing else is used.}

\paragraph{Performance.} Processing \nProcEvents{} events took \wallSec\,s wall time but only
\cpuSec\,s CPU; \picoWallPct\% of the wall time is spent in \fcmd{StPicoDstMaker}, i.e.\ the job is
limited by XRootD I/O, not by the analysis.

\caveat{At this rate (about 30 events/s per job) full-dataset processing needs batch jobs,
ideally scheduled close to the data; this has not been set up yet.}
